\scanpage{79}
% Archive pencil foliation: 78.
de pass\unsure{\deleted{g}a|\deleted{a}g}ge au quotient (sauf erreur de signe, car je n'ai nulle part fait attention aux signes).

8.4. Considérons maintenant l'homomorphisme
\[
(8.4.1)\qquad\begin{tikzcd}[cramped]\handwritten{\underline}{G}(\handwritten{\overline}{\handwritten{\underline}{F}}_p)\arrow[r]&\handwritten{\underline}{G}(\handwritten{\underline}{F}_p)\quad .\end{tikzcd}

\]
Il e\unsure{\deleted{u}s|\deleted{s}u|\deleted{m}s|\deleted{s}m|s}t clair qu'il se factorise par $\handwritten{\underline}{H}$ (i.e. que le composé
\[
\begin{tikzcd}[cramped]\handwritten{\underline}{\unsure{\deleted{i}Q|\deleted{Q}i|\deleted{l}Q|\deleted{Q}l|Q}}_p\text{-faisceaux constructibles sur }\handwritten{\underline}{\unsure{\deleted{P}F|\deleted{F}P|\mathbb{F}|F}}_p\arrow[r,"i",no head,to path={(\tikztostart)--(\tikztotarget)--++(-2pt,2pt)\tikztonodes}]&\operatorname{Fcriso}(\handwritten{\underline}{F}_p)\arrow[r,no head,to path={(\tikztostart)--(\tikztotarget)--++(-2pt,2pt)\tikztonodes}]&\operatorname{Fcriso}(\handwritten{\overline}{\handwritten{\underline}{F}}_p)\end{tikzcd}
\]
s'envoie dans les objets ``constants'' ...) , d'où un homomorphisme
\[
(8.4.2)\qquad\begin{tikzcd}[cramped]\handwritten{\underline}{G}(\handwritten{\overline}{\handwritten{\underline}{F}}_p)\arrow[r,no head,to path={(\tikztostart)--(\tikztotarget)--++(-2pt,2pt)\tikztonodes}]&\handwritten{\underline}{H}\quad .\end{tikzcd}

\]
Par Dieudonné-Manin, on connait la structure des objets de $\operatorname{Fcriso}(\handwritten{\overline}{k})$ pour $\handwritten{\overline}{k}$ algébriquement clos\deleted{e}, on sait que ces objets sont semi-simples et proviennent (mod isomorphisme) d'objets de $\operatorname{Fcriso}(\handwritten{\underline}{F}_p)$. Cela implique que (8.4.1), donc a\unsure{\deleted{a}u|\deleted{u}a|a}s\unsure{\deleted{a}s|\deleted{s}a|a|s}i (8.4.2), est un monomorphisme de gerbes. Pour montrer que c'est une équivalence, il suffit de montrer qu'un objet simple non constant de \deleted{$\handwritten{\underline}{H}$ \unsure{me|ne} donne} $\operatorname{Rep}(\handwritten{\underline}{H})$ donne un objet non constant de \begin{tikzcd}[cramped]\operatorname{Rep}(\handwritten{\underline}{G}(\handwritten{\overline}{k}))\arrow[r,"\handwritten{\approx}",no head,to path={(\tikztostart)--(\tikztotarget)--++(-2pt,2pt)\tikztonodes}]&\operatorname{Fcriso}(\handwritten{\overline}{k}))\end{tikzcd}. Or les \unsure{c\deleted{l}a s\deleted{s}ees|c\deleted{l}a s\deleted{e}ses|c\deleted{a}l s\deleted{s}ees|c\deleted{a}l s\deleted{e}ses|clasees|classes} d'objetssimples de $\operatorname{Rep}(\handwritten{\underline}{H})$ sont indexés par les éléments de $\handwritten{\underline}{Q}$, et si on considère l'objet simple de $\operatorname{Fcriso}(\handwritten{\underline}{F}_p)\handwritten{\simeq}\operatorname{Rep}(\handwritten{\underline}{G}(\handwritten{\underline}{F}_p^{\handwritten{\unsure{\deleted{?}|?|\xi|\deleted{\xi}|\deleted{*}|\deleted{\overline{\vphantom{F}}}}}}\unsure{\rlap{)})|)}$ défini par $(E,\handwritten{\underline}{F}_E)$, avec $\handwritten{\underline}{F}_E$ ayant comme polynôme minimal $T^r-p^s$ avec $r,s\handwritten{\in}\handwritten{\underline}{Z}$, $r\geq1$, $(r,s)=1$, alors sa restriction à $\handwritten{\underline}{H}$ donne un objet simple \deleted{de \unsure{\deleted{k}(|\deleted{(}k|k|(}} (car il reste simple dans $\operatorname{Fcriso}(\handwritten{\overline}{k})$) de classe $s/r\handwritten{\in}\handwritten{\underline}{Q}$, et celui\unsure{\rlap{(}\handwritten{-}|\handwritten{-}}ci, d'après la clas\unsure{\deleted{a}s|\deleted{s}a|s|a}\inserted{\handwritten{i}}ficat\unsure{\deleted{o}i|\deleted{i}o}on de Dieudonné-Manin, \deleted{n'est pas} a\handwritten{\begin{tikzpicture}[baseline=0cm,x=1cm,y=1cm]\draw[thin] (0,-.3)--(.17,.7);\end{tikzpicture}}une image dans $\operatorname{Fc\ iso}(\handwritten{\overline}{k})$ qui n'est pas constante si $r\handwritten{\ne}0$ ...

\handwritten{8\unsure{,|.}5.} On a \unsure{\deleted{e}d|\deleted{d}e}onc \deleted{déterminé} résolu par l'affirmative, pour $k=\handwritten{\underline}{F}_p$, la question 6.2, en déterminant en même temps (à équivalence définie à isomorphisme unique près) la gerbe $\handwritten{\underline}{\unsure{\deleted{G}H|\deleted{H}G|G|H}}(\handwritten{\overline}{k})$ associée à une clôture algébrique $\handwritten{\overline}{k}$ de $k=\handwritten{\underline}{F}_p$ ; \deleted{comme pour} elle ne dépend pas (à \deleted{isomorphi} équivalence déterminée à isomorphisme unique près) du choix de la clôture algébrique envisagée $\handwritten{\overline}{\unsure{\deleted{t}k|\deleted{k}t|k}}$, et par Dieudonné-Manin est can. équi%

